\sscp{The prevalence of synonyms}{02-03-02}
On a technical note before getting further into the discussion
of synonyms, some lexicographers are sceptical that there are very many exact
synonyms (true synonyms, full synonyms),
but accept that there may be many near synonyms
in a language (see \citealp{Zgusta1971}).
We are using the term ‘synonym’ here in a broad sense to refer
to multiple words or phrases for talking about the same or similar things.
Regardless of whether they are exact synonyms,
near synonyms, or only closely related meanings,
the nature of lexical taboos
and parallelisms means that many words that have slightly different
meaning in isolation are treated “as if” they are synonyms when
they are linked together for stylistic and sociolinguistic reasons.

In a development with many parallels to the discussion of PMP
*punti and {\#}muku (both meaning ‘banana’; see \srf{02-03-01-07}),
research summarised in \citet[260]{Schapper2011} notes that a species
of the marsupial cuscus (\it{Phalanger orientalis}) was introduced
by humans westward out from New Guinea into parts of Maluku
and Timor, although the Sulawesi ‘dwarf cuscus’ (\it{Strigocuscus
celenbis}) and the larger ‘bear cuscus’ (\it{Ailurops ursinus}) further
west on Sulawesi are apparently not from human transfer.
All are marsupials of the family \it{Phalangeridae}.
Like the food crops mentioned previously,
this happened several thousand years before the arrival of the
Austronesian speaking peoples into the region.

\citet{Blust1982} reconstructed two terms to
PCEMP, *kandoRa ‘cuscus’ and *mans(aə)r ‘bandicoot’, for marsupials
\ili{east} of the Wallace line.
\citet{Schapper2011} points out that the distribution of glosses
from most languages in our target region actually better supports
PCEMP *mans(aə)r meaning ‘cuscus’,
later only shifting to ‘bandicoot’ in Oceanic languages.
And the evidence supporting the full form of *kandoRa ‘cuscus’
outside of EMP is extremely limited in distribution to only one
language right near the geographical boundary of EMP,
and problematic for several other reasons (see \crf{ch:Mar}).
Reflexes of *kandoRa ‘cuscus’ are found in only a single node
of one of our Wallacean subgroups.

\citet{Blust2012b} acknowledges that Schapper is essentially
correct about the distribution
and glosses of *mans(aə)r as ‘cuscus’ outside of EMP.\footnote{
One of the remaining disputed glosses is discussed further in \srf{02-04-04}.}
He also concedes that “It is,
indeed, puzzling that reflexes of *kandoRa are not more widely
distributed in eastern Indonesia” (\citeyear[267]{Blust2012b}).\footnote{
		We also concede that the apparently conservative form of \ili{Watubela} \it{kadola}
		‘marsupial’ (just west of Indonesian west Papua) is puzzling
		and needs an explanation,
		particularly given that other conservative forms are found far
		away on the eastern end of New Guinea.
		We note that the trisyllabic \ili{Watubela} form is quite different
		in its vowels and structure from formally close words in \tsc{Seram-Tanimbar-Bomberai} and SHWNG languages.
		The comment in \citet{BlustTrussel2020} under the entry for *mansar ‘bandicoot’,
		dismissing \citeauthor{Schapper2011}’s (\citeyear{Schapper2011}) evidence that most languages outside
		of EMP have the gloss of ‘cuscus’ for cognates of *mansar,
		by saying her claim is “based largely on non-specific glosses
		in languages of eastern Indonesia,” lacks integrity.
		We address these issues more fully in \crf{ch:Mar}.}
But he is also dismissive of suggestions to revise his original (\citeyear{Blust1982})
reconstructed glosses and claims of the historical development.
He acknowledges that the two marsupial terms are “a critical
piece of evidence [{\ldots}] for the phylogenetic unity of CEMP”,
although not the sole evidence (\citeyear[262]{Blust2012b}).
The basis of some of his objections reveals his assumptions about
language and the Austronesian dispersal.

From Blust’s perspective, if we accept Schapper’s approach,
we are forced “to accept PCEMP *kandoRa
and *mans(aə)r as synonyms” (\citeyear[267]{Blust2012b}).
Yet, he says, “somewhat awkwardly,
this inference forces her to recognise both *kandoRa
and *mans(aə)r as competitors for the
meaning ‘cuscus’ in PEMP” (\citeyear[267]{Blust2012b}).
And, “we must ask why a language
would have synonyms for ‘cuscus’” (\citeyear[268]{Blust2012b}).

First, his comment is revealing in assuming “a [single] language”.
But the response to his question about synonyms is multi-faceted
and is well documented in the region on all counts.
They include: evidence of many other words with synonyms in the
target region, taboos, parallelisms and doublets, and different patterns of
contact in different parts of the same language.
But first we note that the reconstructed Austronesian vocabulary
in \citet{BlustTrussel2020} has many doublets,
synonyms, and near synonyms at all levels -- as
anything reflecting natural languages should.

